\section{CONCLUSIONES}

El presente proyecto implementó de forma sistemática un ecosistema integral de \textbf{Aprendizaje Automático No Supervisado} para la inteligencia urbana, la criminología computacional y el análisis espaciotemporal, evaluado sobre el benchmark metropolitano de Chicago ($+800.000$ crímenes) y el corpus de noticias de seguridad de Bolivia ($137$ reportes estructurados de los nueve departamentos). A partir de la evidencia analítica y empírica recabada, se establecen las siguientes conclusiones:

\begin{itemize}
    \item \textbf{Superioridad del Clustering Basado en Densidad Geodésica:} La sustitución de distancias euclidianas planas por la métrica de \textit{Haversine} sobre la esfera terrestre demostró ser un requisito matemático indispensable para evitar distorsiones métricas en el análisis geoespacial urbano. Entre los algoritmos evaluados, \textbf{HDBSCAN obtuvo el desempeño más sobresaliente}, alcanzando el mayor Coeficiente de Silueta ($s = 0.435$) y el menor índice de solapamiento de Davies-Bouldin ($0.91$), gracias a su capacidad intrínseca para aislar clústeres de densidades heterogéneas sin requerir un radio $\varepsilon$ uniforme.
    
    \item \textbf{Aislamiento Efectivo de Ruido y Optimización Policial:} DBSCAN y HDBSCAN categorizaron entre el $11.8\%$ y el $14.2\%$ de los incidentes como ruido estocástico (etiqueta $-1$). Este resultado posee un impacto directo en la gestión policial: previene el desplazamiento ineficiente de cuadrillas de patrullaje preventivo hacia eventos aislados que carecen de recurrencia territorial sistemática.
    
    \item \textbf{Potencial de la Minería Semántica No Supervisada (NLP):} La integración de TF-IDF con K-Means ($k=5$) y reducción TruncatedSVD demostró que es posible clasificar narrativas de delitos en lenguaje natural sin requerir etiquetas previas. El sistema descubrió con precisión los arquetipos de \textit{modus operandi} más críticos (asaltos armados en vía pública, hurtos en comercios, violencia doméstica y fraudes financieros), permitiendo habilitar un clasificador en tiempo real con latencia inferior a los $4\text{ ms}$.
    
    \item \textbf{Superación del Problema de la Unidad de Área Modificable (MAUP):} El modelado probabilístico continuo mediante \textbf{Gaussian Mixture Models (GMM)} y la proyección de elipses de covarianza al $95.4\%$ ($2\sigma$) demostraron que el riesgo delictivo se distribuye en superficies continuas y anisotrópicas, superando la rigidez arbitraria de los límites de precintos o distritos policiales tradicionales.
    
    \item \textbf{Detección Neuronal Autónoma de Delitos Anómalos:} La red neuronal Autoencoder con arquitectura de cuello de botella ($D \to 16 \to 4 \to 16 \to D$) demostró una elevada capacidad para aislar delitos excepcionales mediante el error cuadrático medio de reconstrucción (MSE). Fijando el umbral crítico en el percentil 90 ($\tau = 0.0821$), el sistema alertó sobre incidentes infrecuentes de secuestro, homicidios en zonas de nula incidencia histórica y fraudes millonarios, conformando un mecanismo de alerta temprana objetivo.
    
    \item \textbf{Limitaciones Metodológicas Identificadas:} La principal limitación metodológica del proyecto reside en el sesgo de cobertura mediática (\textit{reporting bias}) inherente al corpus de noticias de Bolivia: los medios de comunicación tienden a sobre-representar delitos de alta gravedad (homicidios, feminicidios y atracos armados) mientras que omiten hurtos menores o infracciones contravencionales. Se documenta explícitamente esta restricción para advertir que dicho corpus refleja la agenda mediática delictiva y no necesariamente la totalidad exhaustiva del delito común.
    
    \item \textbf{Transferencia Tecnológica e Impacto Operativo:} El despliegue de la arquitectura web modular en Flask acoplada al visor cartográfico esférico \textbf{Globo 3D en WebGL} demuestra la viabilidad de transferir modelos no supervisados de alta complejidad matemática hacia herramientas interactivas, intuitivas y directamente operables por mandos policiales y gestores de políticas públicas sin formación técnica en ciencia de datos.
\end{itemize}

